\documentclass[11pt,respuestas,a4]{aleph-examen}
% \documentclass[11pt,a5]{aleph-examen}

% -- Paquetes extra
\usepackage{aleph-comandos}
\usepackage{systeme}
\input{formato_ipynb}

% -- Datos
\institucion{Facultad de Ciencias Exactas, Naturales y Ambientales}
\carrera{Catalogo STEM}
\asignatura{Algebra lineal}
\tema{Examen no. 1-1: Matrices}
\autor{Andres Merino}
\fecha{Periodo 2026-1}

\logouno[0.16\textwidth]{Logos/logoPUCE_04_ac}
\definecolor{colortext}{HTML}{1957d1}
\definecolor{colordef}{HTML}{1957d1}
\fuente{montserrat}


\begin{document}

\encabezado

\section*{Indicaciones}
\begin{itemize}[leftmargin=*]
\item
    En esta actividad se evalua si el estudiante (Criterio 2.1) resuelve operaciones con matrices, incluyendo productos y cálculo de inversas; además del cálculo determinantes de matrices y su significado en el contexto del Álgebra Lineal.
\item
    Los ejercicios deben ser resueltos a mano, sin ayuda de resúmenes o asistentes computacionales.
\item
    Las resoluciones de los ejercicios deben ser entregadas en hojas de forma física al docente. No importa el orden de resolucion de los ejercicios.
\item
    Se encuentra prohibido el uso de cualquier otra fuente de información durante todo el examen.
\item
    En caso de considerar que existe un error en la pregunta o que esta se encuentra mal planteada, se debe indicar cual es el error y justificarlo.
\item
    Todas las soluciones deben estar correctamente redactadas y explicadas.
\item
    Se asignarán 4 puntos de excelencia si, además de tener las soluciones correctas, el examen se destaca por su orden y escritura.
\end{itemize}

\section*{Ejercicios}

\begin{preguntas}

%%%%%%%%%%%%%%%%%%%%%%%%%%%%%%%%%%%%%%%%%%
%%%%%%%%%%%%%%%%%%%%%%%%%%%%%%%%%%%%%%%%%%
%%%%%%%%%%%%%%%%%%%%%%%%%%%%%%%%%%%%%%%%%%
\item
    Dada la siguiente matriz
    \[
        A=\begin{pmatrix}
                2 & 0 & 4 \\
                1 & 1 & 0 \\
                -1 & 2 & 3 \\
        \end{pmatrix},
    \]
    realice las siguientes operaciones elementales por filas de forma consecutiva:
    \begin{itemize}
        \item $\frac{1}{2}F_1\rightarrow F_1$ \puntaje{1.5}
        \item $-2F_1+F_2\rightarrow F_2$ \puntaje{1.5}
        \item $-2F_1+F_3\rightarrow F_3$ \puntaje{1.5}
        \item $3F_2+F_1\rightarrow F_2$ \puntaje{1.5}
    \end{itemize}

\begin{respuesta}
    \begin{itemize}
    \item
        $\frac{1}{2}F_1\rightarrow F_1$:
        \[
            \begin{pmatrix}
                1 & 0 & 2 \\
                1 & 1 & 0 \\
                -1 & 2 & 3
            \end{pmatrix}
        \]
    \item
        $-2F_1+F_2\rightarrow F_2$:
        \[
            \begin{pmatrix}
                1 & 0 & 2 \\
                -1 & 1 & -4 \\
                -1 & 2 & 3
            \end{pmatrix}
        \]
    \item
        $-2F_1+F_3\rightarrow F_3$:
        \[
            \begin{pmatrix}
                1 & 0 & 2 \\
                -1 & 1 & -4 \\
                -3 & 2 & -1
            \end{pmatrix}
        \]
    \item
        $3F_2+F_1\rightarrow F_2$: esta no es una operacion elemental por filas valida.\qedhere
    \end{itemize}
\end{respuesta}

%%%%%%%%%%%%%%%%%%%%%%%%%%%%%%%%%%%%%%%%%%
%%%%%%%%%%%%%%%%%%%%%%%%%%%%%%%%%%%%%%%%%%
%%%%%%%%%%%%%%%%%%%%%%%%%%%%%%%%%%%%%%%%%%
\item
    Considere $k\in \R$ y las siguientes matrices:
    \[
        M=\begin{pmatrix}
            k & 1 & 0\\
            0 & 1 & 1
        \end{pmatrix}
        \qquad\text{y}\qquad
        N=\begin{pmatrix}
            1 & 0\\
            0 & 1\\
            1 & 1
        \end{pmatrix},
    \]
    determine para que valores de $k$ el producto $MN$ tiene inversa.\puntaje{11}

\begin{respuesta}
    Primero, calculemos el producto:
    \[
        MN=\begin{pmatrix}
            k & 1\\
            1 & 2
        \end{pmatrix}.
    \]
    La matriz $MN$ tiene inversa si y solo si su determinante es distinto de cero:
    \[
        \det(MN)=\begin{vmatrix}
            k & 1\\
            1 & 2
        \end{vmatrix}
        =2k-1.
    \]
    Por lo tanto, $MN$ es invertible cuando
    \[
        2k-1\neq 0 \iff k\neq \frac{1}{2}.
    \]
    En consecuencia, $MN$ tiene inversa para todo $k\in\R\setminus\left\{\frac{1}{2}\right\}$.\qedhere
\end{respuesta}

%%%%%%%%%%%%%%%%%%%%%%%%%%%%%%%%%%%%%%%%%%
%%%%%%%%%%%%%%%%%%%%%%%%%%%%%%%%%%%%%%%%%%
%%%%%%%%%%%%%%%%%%%%%%%%%%%%%%%%%%%%%%%%%%
\item
    Considere $\alpha\in \R$ y la siguiente matriz:
    \[
        A=\begin{pmatrix}
            1 & \alpha & 0 \\
            0 & 2 & 1 \\
            2 & 1 & 1
        \end{pmatrix},
    \]
    calcula el determinante de $A$ realizando expansion por cofactores a la primera columna.\puntaje{12}

\begin{respuesta}
    Aplicando expansion por cofactores en la primera columna, tenemos:
    \[
        \det(A)=
        1\begin{vmatrix}
            2 & 1 \\
            1 & 1
        \end{vmatrix}
        -0\begin{vmatrix}
            \alpha & 0 \\
            1 & 1
        \end{vmatrix}
        +2\begin{vmatrix}
            \alpha & 0 \\
            2 & 1
        \end{vmatrix}.
    \]
    Entonces:
    \[
        \det(A)=1(2\cdot 1-1\cdot 1)+2(\alpha\cdot 1-0\cdot 2)
        =1+2\alpha.
    \]
    Por lo tanto,
    \[
        \det(A)=1+2\alpha.\qedhere
    \]
\end{respuesta}

%%%%%%%%%%%%%%%%%%%%%%%%%%%%%%%%%%%%%%%%%%
%%%%%%%%%%%%%%%%%%%%%%%%%%%%%%%%%%%%%%%%%%
%%%%%%%%%%%%%%%%%%%%%%%%%%%%%%%%%%%%%%%%%%
\item
    Suponga que la siguiente matriz es no singular, calcule, usando reduccion por filas, la inversa:\puntaje{12}
    \[
        A=\begin{pmatrix}
            1 & 1 & 0 \\
            0 & 2 & 2 \\
            0 & 1 & 3
        \end{pmatrix}.
    \]

\begin{respuesta}
    Partimos de
    \[
        \begin{pmatrix}
            1 & 1 & 0 & | & 1 & 0 & 0 \\
            0 & 2 & 2 & | & 0 & 1 & 0 \\
            0 & 1 & 3 & | & 0 & 0 & 1
        \end{pmatrix},
    \]
    y aplicamos operaciones elementales:
    \begin{align*}
        \begin{pmatrix}
            1 & 1 & 0 & | & 1 & 0 & 0 \\
            0 & 2 & 2 & | & 0 & 1 & 0 \\
            0 & 1 & 3 & | & 0 & 0 & 1
        \end{pmatrix}
        &\sim
        \begin{pmatrix}
            1 & 1 & 0 & | & 1 & 0 & 0 \\
            0 & 1 & 3 & | & 0 & 0 & 1 \\
            0 & 2 & 2 & | & 0 & 1 & 0
        \end{pmatrix} && F_2\leftrightarrow F_3 \\
        &\sim
        \begin{pmatrix}
            1 & 1 & 0 & | & 1 & 0 & 0 \\
            0 & 1 & 3 & | & 0 & 0 & 1 \\
            0 & 0 & -4 & | & 0 & 1 & -2
        \end{pmatrix} && -2F_2+F_3\rightarrow F_3 \\
        &\sim
        \begin{pmatrix}
            1 & 1 & 0 & | & 1 & 0 & 0 \\
            0 & 1 & 3 & | & 0 & 0 & 1 \\
            0 & 0 & 1 & | & 0 & -\frac{1}{4} & \frac{1}{2}
        \end{pmatrix} && -\frac{1}{4}F_3\rightarrow F_3 \\
        &\sim
        \begin{pmatrix}
            1 & 1 & 0 & | & 1 & 0 & 0 \\
            0 & 1 & 0 & | & 0 & \frac{3}{4} & -\frac{1}{2} \\
            0 & 0 & 1 & | & 0 & -\frac{1}{4} & \frac{1}{2}
        \end{pmatrix} && -3F_3+F_2\rightarrow F_2 \\
        &\sim
        \begin{pmatrix}
            1 & 0 & 0 & | & 1 & -\frac{3}{4} & \frac{1}{2} \\
            0 & 1 & 0 & | & 0 & \frac{3}{4} & -\frac{1}{2} \\
            0 & 0 & 1 & | & 0 & -\frac{1}{4} & \frac{1}{2}
        \end{pmatrix}
        && -F_2+F_1\rightarrow F_1.
    \end{align*}
    Por lo tanto,
    \[
        A^{-1}=\begin{pmatrix}
            1 & -\frac{3}{4} & \frac{1}{2} \\
            0 & \frac{3}{4} & -\frac{1}{2} \\
            0 & -\frac{1}{4} & \frac{1}{2}
        \end{pmatrix}.\qedhere
    \]
\end{respuesta}

\end{preguntas}

\end{document}
